\chapter{Dystonia}

\section*{Definition and Overview}
Dystonia is a movement disorder in which muscles contract involuntarily, causing twisting, repetitive movements or abnormal postures. It can affect a single muscle group, such as the neck, eyes, or hand, or involve many muscles throughout the body. Dystonia is caused by a problem in the part of the brain that controls movement, and it is not a sign of weakness, illness, or mental health problems. There are many forms of dystonia, and treatment, including medications and botulinum toxin injections, is very effective for many people.

\section*{Epidemiology}
Dystonia is the third most common movement disorder after Parkinson's disease and tremor. Focal dystonias, which affect one part of the body, are the most common form and include cervical dystonia (torticollis), blepharospasm (involuntary eye closure), and writer's cramp. Dystonia can begin at any age, from childhood to late adulthood, and affects men and women across all populations. Many cases are mild and never formally diagnosed.

\section*{Aetiology and Pathophysiology}
Dystonia arises from dysfunction in the basal ganglia, the deep brain structures that help control and coordinate movement.
\begin{itemize}
    \item \textbf{Primary dystonia:} no identifiable cause other than a genetic predisposition; many genes have been identified, including DYT1
    \item \textbf{Secondary dystonia:} caused by brain injury, stroke, infection, or exposure to certain drugs
    \item \textbf{Drug-induced dystonia:} antipsychotic medicines and some other drugs can cause acute or persistent dystonia
    \item \textbf{Mechanism:} abnormal signals from the basal ganglia cause excessive, sustained muscle contraction and impaired inhibition of unwanted movement
    \item \textbf{Age of onset:} childhood-onset dystonia often becomes generalised, while adult-onset dystonia tends to remain focal
    \item \textbf{Triggers and patterns:} symptoms may be task-specific, appearing only during actions such as writing or speaking
\end{itemize}

\section*{Clinical Features}
The features depend on which muscles are affected.
\begin{itemize}
    \item Involuntary twisting or repetitive movements of a body part
    \item Abnormal postures, such as the head pulling to one side, tilting, or turning (cervical dystonia)
    \item Rapid, involuntary blinking or forced eye closure (blepharospasm)
    \item Cramping of the hand during writing (writer's cramp)
    \item A strained, strangled voice (spasmodic dysphonia)
    \item Tremor, pain, and fatigue in the affected muscles
    \item Symptoms that may worsen with stress, fatigue, or specific tasks
    \item In generalised dystonia: involvement of the trunk, limbs, and whole body
\end{itemize}

\section*{Differential Diagnosis}
Other movement and neurological conditions must be considered.
\begin{itemize}
    \item Parkinson's disease and other parkinsonian disorders
    \item Essential tremor, which causes rhythmic shaking rather than twisting
    \item Tics and Tourette syndrome
    \item Cerebral palsy and other static brain injuries
    \item Functional (psychogenic) movement disorders
    \item Stroke, multiple sclerosis, and other causes of spasticity
\end{itemize}

\section*{When to Seek Medical Care}
Anyone with involuntary muscle contractions, twisting movements, or abnormal postures should see a GP and be referred to a neurologist. Sudden onset of dystonia, or dystonia after starting a new medication, requires prompt review.

\textbf{Call 000 (or local emergency services) for:}
\begin{itemize}
    \item Sudden severe neck spasm with difficulty breathing or swallowing
    \item Sudden weakness, facial drooping, or difficulty speaking, which may indicate stroke
    \item Severe acute reactions to medication, including uncontrolled muscle spasms
    \item Collapse or loss of consciousness
\end{itemize}

\section*{Diagnosis and Investigations}
Diagnosis is largely clinical, based on the pattern of movement.
\begin{itemize}
    \item \textbf{History and examination:} the nature, distribution, and triggers of the movements, and any medications
    \item \textbf{Neurological examination:} assessment of the movements and exclusion of other movement disorders
    \item \textbf{Genetic testing:} for specific forms of inherited dystonia, particularly in early-onset or familial cases
    \item \textbf{Imaging:} MRI of the brain to exclude structural causes
    \item \textbf{Blood tests:} including tests for Wilson's disease and other metabolic causes in selected cases
    \item \textbf{Medication review:} identification of any drugs that can cause or worsen dystonia
\end{itemize}

\section*{Management}
Treatment is tailored to the type and severity of dystonia.
\begin{itemize}
    \item \textbf{Botulinum toxin injections:} the first-line treatment for most focal dystonias; the toxin weakens the overactive muscles, relieving spasm for several months
    \item \textbf{Oral medications:} muscle relaxants, anticholinergics, and other medicines help some people, particularly with generalised dystonia
    \item \textbf{Deep brain stimulation:} a surgical treatment in which electrodes are implanted in the brain to regulate abnormal signals; very effective for many people with generalised or cervical dystonia
    \item \textbf{Physiotherapy and occupational therapy:} stretching, splinting, and strategies to manage daily tasks
    \item \textbf{Speech therapy:} for spasmodic dysphonia
    \item \textbf{Treatment of the cause:} adjustment of offending medications and treatment of underlying conditions
\end{itemize}

\section*{Complications}
\begin{itemize}
    \item Pain and fatigue in the affected muscles
    \item Reduced ability to work, drive, and perform daily tasks
    \item Social embarrassment and withdrawal
    \item Depression and anxiety
    \item In severe generalised dystonia: difficulty walking, eating, and breathing
    \item Complications of treatment, including side effects of medicines and surgery
\end{itemize}

\section*{Prognosis and Outcomes}
The outlook for dystonia varies. Focal dystonias respond well to botulinum toxin, and most people achieve good symptom control. Generalised dystonia is more challenging but often improves substantially with deep brain stimulation. Dystonia is not progressive in most forms and does not shorten life, but it is usually a lifelong condition that requires ongoing management.

\section*{Prevention and Self-Care}
\begin{itemize}
    \item Seek specialist neurological assessment for an accurate diagnosis
    \item Learn to recognise and avoid triggers such as stress and fatigue
    \item Use sensory tricks that may temporarily relieve symptoms, such as touching the chin in cervical dystonia
    \item Maintain strength and flexibility with physiotherapy
    \item Manage stress with relaxation techniques
    \item Connect with support groups for information and community
    \item Take medications as prescribed and report side effects
\end{itemize}

\section*{Sources and Further Reading}
The following reputable sources provide further information for healthcare students and patients seeking reliable, current advice about dystonia.
\begin{itemize}
    \item Dystonia Australia. \textit{Support and information for people with dystonia.} https://www.dystonia.org.au/
    \item NHS. \textit{Dystonia: overview and treatment.} https://www.nhs.uk/conditions/dystonia/
    \item Mayo Clinic. \textit{Dystonia: symptoms and causes.} https://www.mayoclinic.org/diseases-conditions/dystonia/
    \item MedlinePlus. \textit{Dystonia.} https://medlineplus.gov/dystonia.html
    \item Dystonia Medical Research Foundation. \textit{Patient education and resources.} https://dystonia-foundation.org/
    \item MSD Manual. \textit{Dystonia: professional overview.} https://www.msdmanuals.com/
\end{itemize}
